% !PRINTED 710
\weberpage{710}
%<<CONT
elementaren Sätzen der Integralrechnung, \gesperrt{wenn $\alpha$ irgend einen Winkel
zwischen 0 und $2\pi$ bedeutet,}
\[
\int_0^1 \frac{dx}{x-e^{i\alpha}} = \tfrac{1}{2}\log\left(4\sin^2\frac{\alpha}{2}\right) + i\,\frac{\pi-\alpha}{2}.
\]
Hieraus ergiebt sich nach (4):
\weq{5}{X = -\frac{1}{m}\underset{1,\,m-1}{\overset{s}{\Sigma}} f(r^s)\left[\tfrac{1}{2}\log 4\left(\sin\frac{s\pi}{m}\right)^2 - \frac{i\pi(m-2s)}{2m}\right].}
Hierin ist nach (2)
\weq{6}{f(r^s) = \overset{t}{\Sigma}\chi(t)\,r^{ts},}
und daraus ergiebt sich [wegen (3)]:
\weq{7}{\underset{1,\,m-1}{\overset{s}{\Sigma}} f(r^s) = \underset{0,\,m-1}{\overset{s}{\Sigma}} f(r^s) = 0,}
da
\[
\underset{0,\,m-1}{\overset{s}{\Sigma}} r^{st} = 0
\]
ist. Hiernach vereinfacht sich die Formel (5) noch etwas und ergiebt
\weq{8}{X = -\frac{1}{m}\underset{1,\,m-1}{\overset{s}{\Sigma}} f(r^s)\left[\tfrac{1}{2}\log\left(\sin\frac{s\pi}{m}\right)^2 + \frac{i\pi s}{m}\right],}
wofür man auch, wenn man $s$ durch $m-s$ ersetzt und wieder (7) benutzt,
\weq{9}{X = -\frac{1}{m}\underset{1,\,m-1}{\overset{s}{\Sigma}} f(r^{-s})\left[\tfrac{1}{2}\log\left(\sin\frac{s\pi}{m}\right)^2 - \frac{i\pi s}{m}\right]}
setzen kann.

Aus (6) folgt aber
\[
f(r^{-s}) = \overset{t}{\Sigma}\chi(t)\,r^{-ts},
\]
und die nach $t$ genommene Summe ändert sich nicht, wenn $t$ durch $m-t$
ersetzt wird. Daraus folgt:
\[
f(r^{-s}) = \chi(-1)\,f(r^s).
\]
Der Factor $\chi(-1)$ hat den Werth $+1$ oder $-1$, da sein Quadrat $=+1$
ist. Wir unterscheiden daher jetzt zwei Arten von Charakteren $\chi_1(n)$,
$\chi_2(n)$, so dass
\weq{10}{\chi_1(-1) = 1,\quad \chi_2(-1) = -1}
%CONT>>
